%%%%%%%%%%%%%%%%%%%%%%%%%%%%%%%%%%%%%%%%%%%%%%%%%%%%%%%%%%%%%%%%%%%%%%%%%%%
% Función: estructura del documento.
%
% Uso: copia de restauración; normalmente no debes modificarla.
% Tras restaurarla, modifica la copia activa para adaptarla a tu trabajo.
%%%%%%%%%%%%%%%%%%%%%%%%%%%%%%%%%%%%%%%%%%%%%%%%%%%%%%%%%%%%%%%%%%%%%%%%%%%

\chapter{Estructura del documento}
\label{cha:estructura-documento}

En este capítulo recorremos las partes que forman el documento principal para que sepas dónde escribir cada contenido. Comenzaremos por su organización general, seguiremos el orden de las páginas iniciales, el cuerpo y el cierre, y terminaremos distinguiendo qué partes tendrás que modificar. Las describimos todas para que tengas una referencia completa, pero no tendrás que modificar muchas de ellas porque la plantilla las genera o selecciona automáticamente.

\section{Organización general del documento}
\label{sec:estr-del-docum}

\texttt{Book/book.tex} es el punto de entrada estable que compilas, mientras que \texttt{Book/content-standard.tex} organiza los capítulos, la bibliografía y los apéndices de los documentos habituales. En este apartado describimos la organización general para que puedas localizar cada bloque sin tener que modificar la infraestructura de \texttt{book.tex}.

El documento se divide en tres partes gestionadas por las órdenes \texttt{\textbackslash{}frontmatter}, \texttt{\textbackslash{}mainmatter} y \texttt{\textbackslash{}backmatter}:

\begin{itemize}
  \item \texttt{\textbackslash{}frontmatter}: contiene las cubiertas, los documentos preliminares, la dedicatoria, los agradecimientos, los resúmenes y los distintos índices y listados.
  \item \texttt{\textbackslash{}mainmatter}: contiene los capítulos principales, los apartados técnicos opcionales, la bibliografía y los apéndices.
  \item \texttt{\textbackslash{}backmatter}: contiene la contraportada cuando corresponde al tipo de trabajo seleccionado.
\end{itemize}

Estas órdenes controlan la numeración y la organización interna del libro. No necesitas modificarlas ni cambiar su posición.

\section{Páginas iniciales}
\label{sec:paginas-iniciales}

La parte inicial del documento, delimitada mediante \texttt{\textbackslash{}frontmatter}, reúne las cubiertas, el contenido preliminar y los índices. La plantilla distribuida incluye los siguientes elementos:

\begin{itemize}
  \item Portada y páginas iniciales, seleccionadas automáticamente por \texttt{cover/cover.tex} a partir de la titulación configurada mediante \texttt{\textbackslash{}myDegree}. Normalmente solo tendrás que configurar los datos de \texttt{Config/myconfig.tex}, sin modificar \texttt{cover/cover.tex} ni las plantillas institucionales, organizadas por universidad en \texttt{cover/uah/}, \texttt{cover/upm/} y \texttt{cover/urjc/}.
  \item Documentos PDF preliminares, opcionales y desactivados de forma predeterminada. El bloque de \texttt{\textbackslash{}includepdf} en \texttt{Book/book.tex} permite incorporar una carta de autorización u otro PDF antes de la dedicatoria. Para utilizarlo, sustituye allí la ruta del ejemplo \texttt{letters/sampleLetter.pdf} por la de tu documento y activa \texttt{\textbackslash{}myIncludeSampleLetter} en \texttt{Config/myconfig.tex}; \texttt{pages=-} incluye todas sus páginas.
  \item Dedicatoria y agradecimientos, incluidos desde \texttt{dedication/dedicatoria.tex} y \texttt{acknowledgements/agradecimientos.tex}. Puedes mostrarlos u ocultarlos mediante \texttt{\textbackslash{}myIncludeDedication} y \texttt{\textbackslash{}myIncludeAcknowledgements}.
  \item Definiciones de acrónimos y símbolos, incluidas desde \texttt{acronyms/defacronymsgl.tex} y \texttt{symbols/defsymbolsgl.tex}. Sus listas se controlan mediante \texttt{\textbackslash{}myIncludeListOfAcronyms} y \texttt{\textbackslash{}myIncludeListOfSymbols}; las secciones~\ref{sec:uso-de-acronimos} y~\ref{sec:simbolos} explican su definición y uso.
  \item Resumen en español, incluido desde \texttt{abstract/resumen.tex}, y resumen en inglés, incluido desde \texttt{abstract/abstract.tex}. Normalmente tendrás que editar ambos ficheros con el contenido de tu trabajo.
  \item Resumen extendido en español, disponible en \texttt{abstract/resumen-extendido.tex}. Su inclusión está comentada de forma predeterminada; actívala únicamente cuando lo requiera tu tipo de trabajo.
  \item Índice general, incluido siempre desde \texttt{cover/table-of-contents.tex}, e índices de figuras y tablas, generados respectivamente desde \texttt{cover/list-of-figures.tex} y \texttt{cover/list-of-tables.tex}. Estos dos últimos se controlan mediante \texttt{\textbackslash{}myIncludeListOfFigures} y \texttt{\textbackslash{}myIncludeListOfTables}.
  \item Listados adicionales de código, algoritmos y vídeos, generados desde \texttt{cover/list-of-code.tex}, \texttt{cover/list-of-algorithms.tex} y \texttt{cover/list-of-videos.tex}, y controlados mediante las variables \texttt{\textbackslash{}myIncludeListOfCode}, \texttt{\textbackslash{}myIncludeListOfAlgorithms} y \texttt{\textbackslash{}myIncludeListOfVideos}. La sección~\ref{sec:intro-videolink} explica el ejemplo dedicado a los vídeos.
\end{itemize}

\section{Cuerpo principal}
\label{sec:cuerpo-principal}

La orden \texttt{\textbackslash{}mainmatter} da paso al contenido principal y activa la numeración ordinaria de los capítulos. En la plantilla encontrarás estos bloques:

\begin{itemize}
  \item Capítulos principales, incluidos explícitamente desde el directorio \texttt{chapters/} mediante las órdenes \texttt{\textbackslash{}input} de \texttt{Book/content-standard.tex}. El documento distribuido contiene esta guía, un capítulo de elementos básicos que incluye ejemplos de figuras y tablas, y otro de composición avanzada. Cuando prepares tu trabajo, sustituye esta selección por tus propios capítulos o utiliza la estructura mínima descrita en la sección~\ref{sec:estructura-minima}.
  \item Pliego de condiciones, disponible en \texttt{pliego/pliego.tex}, y presupuesto, disponible en \texttt{presupuesto/presupuesto.tex}. Sus órdenes de inclusión están comentadas de forma predeterminada; actívalas únicamente cuando necesites estos documentos. Los ejemplos proceden de un trabajo fin de carrera de Jesús Martínez y se conservan como material de referencia; debido a su antigüedad, contrástalos con la normativa y los formatos actuales.
  \item Bibliografía, impresa mediante \texttt{biblio/bibliography.tex}. Añade tus referencias a los ficheros \texttt{.bib} declarados en \texttt{biblio/bibliofiles.tex}; normalmente no necesitarás modificar \texttt{biblio/bibliography.tex}.
  \item Apéndices, agrupados mediante el entorno \texttt{appendices} e incluidos explícitamente mediante las órdenes \texttt{\textbackslash{}input} de \texttt{Book/content-standard.tex}. El documento distribuido dedica un apéndice al código fuente, los algoritmos y los listados adicionales. Puedes reemplazarlos, añadir otros nuevos u omitirlos según las necesidades de tu trabajo.
\end{itemize}

\section{Estructura especializada para tesis por compendio}

Solo si preparas una tesis doctoral por compendio de publicaciones debes seleccionar \texttt{compendium} mediante \texttt{\textbackslash{}myDocumentStructure} y organizar su cuerpo en \texttt{Book/content-compendium.tex}. Este fichero sustituye la selección convencional de capítulos por un resumen extendido, sitúa su bibliografía antes de la segunda parte y declara los PDF de las publicaciones. La sección~\ref{sec:tesis-compendio} contiene las instrucciones completas; el resto de usuarios puede ignorar esta estructura especializada.

\section{Cierre del documento}
\label{sec:cierre-documento}

La orden \texttt{\textbackslash{}backmatter} delimita el cierre del documento. En la estructura actual da paso a la contraportada, que se selecciona automáticamente mediante \texttt{cover/backpage.tex} cuando existe una plantilla aplicable al tipo de trabajo y a la titulación configurados. Normalmente no necesitarás modificar este fichero ni la plantilla institucional seleccionada.

\section{Qué partes tendrás que modificar}
\label{sec:partes-modificar}

Aunque hemos descrito todos los componentes para que puedas utilizar este capítulo como referencia, en la práctica solo tendrás que intervenir en una parte de ellos. Puedes distinguirlos de la siguiente manera:

\begin{itemize}
  \item Contenido que normalmente debes editar: los capítulos, los resúmenes y la bibliografía.
  \item Contenido que debes editar o activar únicamente si lo utilizas: la dedicatoria, los agradecimientos, los PDF preliminares, el resumen extendido, los acrónimos, los símbolos, los listados adicionales, el pliego de condiciones, el presupuesto y los apéndices.
  \item Elementos que normalmente no debes modificar: las órdenes estructurales, la generación de los índices y la selección automática de las cubiertas y contraportadas.
\end{itemize}

Modifica las órdenes de inclusión de \texttt{Book/content-standard.tex} cuando necesites añadir, omitir o reordenar capítulos y apéndices, y utiliza las variables \texttt{\textbackslash{}myInclude...} descritas en la sección~\ref{sec:elementos-opcionales-configuracion} para el contenido preliminar y los listados. Normalmente no debes modificar \texttt{Book/book.tex}. Si preparas la modalidad doctoral especializada, realiza en cambio la organización del cuerpo y de las publicaciones en \texttt{Book/content-compendium.tex}, siguiendo la sección~\ref{sec:tesis-compendio}.


\section{Resumen del capítulo}
\label{sec:resumen-estructura-documento}

El documento se organiza mediante \texttt{\textbackslash{}frontmatter}, \texttt{\textbackslash{}mainmatter} y \texttt{\textbackslash{}backmatter}. En la práctica trabajarás principalmente con \texttt{Config/myconfig.tex}, \texttt{Book/content-standard.tex}, los capítulos, los resúmenes y la bibliografía; activarás otros contenidos mediante las variables \texttt{\textbackslash{}myInclude...} y dejarás intactos \texttt{Book/book.tex} y los mecanismos automáticos de cubiertas. La excepción es la tesis doctoral por compendio, cuyo cuerpo se organiza mediante \texttt{Book/content-compendium.tex}.

%%% Local Variables:
%%% TeX-master: "../book"
%%% End:
